% ============================================================
% abstract.tex
% Tesis: Impacto del Gasto en Limpieza Pública sobre
%        Gestión de Residuos Sólidos Municipales, Perú 2015-2019
% Autor: Dante Barreto Gamarra (UNAS)
% Nota: 247 palabras --- excede INV-5 (150 palabras) por instrucción
%       explícita del autor; reducir antes de envío a revista.
%       Todas las citas metodológicas usan claves de Bibliography_base.bib.
% ============================================================

\begin{abstract}
\noindent \singlespacing
Este estudio estima el efecto del gasto municipal en limpieza pública sobre la gestión de residuos sólidos en los distritos del Perú durante el período 2015--2019, en el marco de la Ley Orgánica de Municipalidades (Ley N° 27972) y el proceso de modernización del Estado. La principal amenaza de identificación es la endogeneidad del gasto: los municipios con mayor capacidad institucional ejecutan más presupuesto y reportan mejores indicadores simultáneamente. Para corregir este sesgo, se aplica el estimador de Efectos Aleatorios Correlacionados con función de control (CRE+CF) propuesto por \citet{Wooldridge2015}, que combina el dispositivo de \citet{Mundlak1978} con residuos de primera etapa como regresor adicional en los modelos de segunda etapa. El análisis cubre tres dimensiones de gestión: eficiencia en recolección, planeamiento municipal y disposición final adecuada. Los resultados, obtenidos sobre un panel de 9,322 observaciones, 1,874 municipios y 196 clusters provinciales, muestran que un incremento del 10\% en el gasto devengado eleva en 8{,}79\% la cantidad diaria de residuos recolectados (elasticidad $= +0{,}879$; $p < 0{,}001$), aumenta el índice de planeamiento municipal en 0{,}030 puntos sobre una escala de cinco instrumentos ($\hat{\beta} = +0{,}300$; $p < 0{,}001$), e incrementa en 2{,}2 puntos porcentuales la probabilidad de disposición formal en relleno sanitario sobre una línea base del $18\%$ de municipios (AME $= +0{,}022$; $p < 0{,}001$). Las tres hipótesis primarias son robustas a la corrección Bonferroni-Holm. Los resultados identifican al financiamiento municipal como la restricción vinculante para la formalización ambiental de la disposición final, y al gasto como habilitador de la capacidad institucional ---medida por la adopción del PIGARS--- que la regulación peruana vigente (Decreto Legislativo N° 1278) exige como precondición para el acceso a infraestructura de disposición. La evidencia respalda reformas en la fórmula del FONCOMUN orientadas a municipios con mayor déficit de gestión de residuos sólidos.
\end{abstract}

\vspace{1em}
\noindent \textbf{Códigos JEL:} H72, Q53, R51, O18, H77

\vspace{0.5em}
\noindent \textbf{Palabras clave:} gasto municipal, residuos sólidos, relleno sanitario, CRE+CF, Mundlak, Perú, FONCOMUN, DL 1278, Ley 27972
